% TOP_PROBLEM: {"id":30007615,"problem_number":"LOCAL-30007615","title":"Global financial-cycle transmission","category_id":21,"category":{"id":21,"name":"economics","display_name":"Economics","description":"Open economic theory statements, published research agendas, and proposed scoped specifications; question provenance and historical priority are qualified per record.","slug":"economics","order_index":21,"created_at":"2026-10-08T00:00:00Z"},"status":"research_question","external_url":null,"legacy_ids":[],"published":false,"statement":"In the explicitly specified setting: How do safe-asset demand, dollar funding and nonbank balance sheets transmit global financial shocks to different economies? The mathematical statement above fixes the target, assumptions and scope of this catalogue formulation.","economics":{"original_id":104,"field":"International finance and exchange rates","kind":"identification","formalization_basis":"adapted_research_question","source_status":"research_agenda","priority_status":"earliest_found","priority_note":"This is the earliest explicit statement verified in this audit; first priority for the full catalogue formulation is not established.","earliest_verified_reference":{"authors":["Bank of England"],"title":"Bank of England Agenda for Research, 2025–2028","year":2026,"url":"https://www.bankofengland.co.uk/research/bank-of-england-agenda-for-research","doi":null,"locator":"Interconnected World / Global drivers, paragraph 3; Priority topics for 2026, item 3","evidence_summary":"Requests global-financial-cycle channels and intermediary transmission research."},"current_reference":{"authors":["Bank of England"],"title":"Bank of England Agenda for Research, 2025–2028","year":2026,"url":"https://www.bankofengland.co.uk/research/bank-of-england-agenda-for-research","doi":null,"locator":"Interconnected World / Global drivers, paragraph 3; Priority topics for 2026, item 3","evidence_summary":"Requests global-financial-cycle channels and intermediary transmission research."},"formalization_note":"Adapts an explicitly verified research-agenda passage. Mathematical objects, interventions, objectives and scope are proposed here; existing model-specific solutions are not relabeled unsolved."},"statusQualification":"Published question or research agenda verified at the cited locator. The mathematical specification is adapted for this catalogue and includes additional assumptions; source publication does not establish first-ever priority or certify this exact model remains unsolved. Adapts an explicitly verified research-agenda passage. Mathematical objects, interventions, objectives and scope are proposed here; existing model-specific solutions are not relabeled unsolved.","statusReviewedAt":"2026-10-08"}
% FIRST_POSTED: 2026-10-08
% ENTRY_KIND: statement_only
% !TeX program = lualatex
\documentclass[11pt]{article}
\usepackage{fontspec}
\usepackage[a4paper,margin=25mm]{geometry}
\usepackage{amsmath,amssymb,mathtools}
\usepackage{xurl}
\usepackage[hidelinks]{hyperref}
\newcommand{\sref}[2]{\hyperlink{source:#1:#2}{[S#2]}}
\setlength{\emergencystretch}{3em}
\begin{document}
% problemId:30007615
\section*{Global financial-cycle transmission}\label{30007615}
\subsection{Definitions and mathematical statement}
Fix countries \(c=1,\ldots,C\), intermediary balance sheets, dollar funding exposures and demand for safe dollar claims. Let \(z_t\) be an externally identified global safe-asset-demand shock, normalized to unit variance. Let \(Y_{c,t+h}\) denote a declared outcome such as output, local credit or asset prices. Define total responses \(\tau_c(h)=\partial E[Y_{c,t+h}(do(z_t=u))]/\partial u|_{u=0}\). To separate funding and nonbank channels, specify counterfactual regimes that hold each channel's transmission equation at its pre-shock value while allowing all remaining equations to adjust. The task is to identify total and channel-specific responses, including interactions that prevent contributions from adding mechanically. Establish the shock's exclusion from unrelated productivity or policy news. Distinguish common shocks from bilateral spillovers and identify which dollar exposures create amplification rather than merely reflecting prior risk. Compare countries using prespecified initial states and horizon sets. Supply bounds when channel interventions are not separately observed. Do not equate principal components of global asset prices with an identified causal shock. Report whether the resulting decomposition is stable across crisis and ordinary conditions and across bank-dominated versus market-based financial systems.

\par\textbf{Formulation and scope.} Adapts an explicitly verified research-agenda passage. Mathematical objects, interventions, objectives and scope are proposed here; existing model-specific solutions are not relabeled unsolved.

\par\textbf{Open-question provenance.} An explicit question or research agenda was verified at the source locator. This does not by itself certify that every later version remains unresolved \sref{30007615}{1}.

\par\textbf{Historical priority.} This is the earliest explicit statement verified in this audit; first priority for the full catalogue formulation is not established.

\subsection{Short English statement}
In the explicitly specified setting: How do safe-asset demand, dollar funding and nonbank balance sheets transmit global financial shocks to different economies? The mathematical statement above fixes the target, assumptions and scope of this catalogue formulation.

\subsection{Sources}
\begin{itemize}\raggedright
\item[\textnormal{[S1]}]\hypertarget{source:30007615:1}{} Bank of England. 2026. Bank of England Agenda for Research, 2025–2028. Interconnected World / Global drivers, paragraph 3; Priority topics for 2026, item 3.
\url{https://www.bankofengland.co.uk/research/bank-of-england-agenda-for-research}.
{\small\textit{Earliest verified question/agenda reference; historical priority qualified below. Requests global-financial-cycle channels and intermediary transmission research.}}
\end{itemize}
\end{document}
